% Memberikan label pada bab ini
\renewcommand{\thechapter}{\Roman{chapter}}
\chapter{PENUTUP}\label{babLima}
\renewcommand{\thechapter}{\arabic{chapter}}
\vspace{8mm}

% Penjelasan umum mengenai isi bab
% {\frenchspacing
%     Berdasarkan rumusan masalah dan hasil perhitungan algoritma \textit{particle swarm optimization} (PSO) untuk optimalisasi pemilihan urutan \textit{scene} syuting film "Ketika Adzan Sudah Tidak Lagi Berkumandang"
%     dapat diambil kesimpulan sebagai berikut:
% }

% Sub Bab
\section{Kesimpulan}
\vspace{-4mm}
{\frenchspacing
    Kesimpulan yang dapat diambil dari hasil penelitian Optimalisasi Pemilihan Urutan Syuting \textit{Scene} Film Menggunakan Algoritma \textit{Particle Swarm Optimization} (PSO) (Studi Kasus: Film "Ketika Adzan Sudah Tidak Lagi Berkumandang") adalah sebagai berikut:

    \begin{enumerate}[align = left, left = 0cm, nolistsep]
        \item Model graf dari pemilihan urutan syuting \textit{scene} film "Ketika Adzan Sudah Tidak Lagi Berkumandang" merupakan graf lengkap yang dapat dilihat pada Gambar \ref{gam: Graf Model}.
              Tiap \textit{scene} dalam data \textit{scene} dianalogikan sebagai \textit{vertex} dan bobot biaya syuting antar \textit{scene} sebagai \textit{edge}.
              Model matematika pada pemilihan urutan syuting \textit{scene} film "Ketika Adzan Sudah Tidak Lagi Berkumandang" dibentuk menggunakan konsep \textit{traveling salesman problem} (TSP) adalah sebagai berikut.

              \noindent
              Fungsi tujuan
              \begin{equation}
                  \label{pers: model matematika penelitian}
                  \operatorname{min} f(a_{j,k}) = \sum^{59}_{i=0} \sum^{59}_{j=0} a_{j,k}(v_{j}v_{k})
              \end{equation}

              \noindent
              Fungsi kendala
              \begin{equation*}
                  \begin{aligned}
                       & \sum^{59}_{j=0} a_{j,k} = 1,    & k = 0,1,\dots,59 \\
                       & \sum^{59}_{k=0} a_{j,k} = 1,    & j = 0,1,\dots,59 \\
                       & a_{j,k} \in \left\{0,1\right\}, & j \neq k
                  \end{aligned}
              \end{equation*}

              \noindent
              Dimana $v_{j}v_{k}$ merupakan elemen matriks biaya antar \textit{scene} dari \textit{scene} ke-$j$ menuju \textit{scene} ke-$k$ sesuai Tabel \ref{tab: data biaya antar scene}
              dan $a_{j,k}$ merupakan elemen matriks biner solusi dari TSP seperti pada Tabel \ref{tab: contoh matriks aij}.

        \item Algoritma \textit{particle swarm optimization} (PSO) untuk optimalisasi pemilihan urutan syuting \textit{scene} film "Ketika Adzan Sudah Tidak Lagi Berkumandang" yang diterapkan meliputi proses normalisasi data biaya antar \textit{scene}, proses algoritma PSO, pengujian modifikasi parameter PSO, dan analisis hasil perhitungan algoritma PSO.
              Proses perhitungan algoritma PSO menggunakan nilai parameter terbaik yang diperoleh dari hasil pengujian modifikasi nilai parameter PSO.
              Nilai parameter terbaik yang dipakai yaitu iterasi maksimum $(n)$ sebesar $900$, ukuran \textit{swarm} $(N)$ sebesar $250$, dan faktor \textit{clamping} $(\varepsilon)$ pada \textit{velocity} maksimum sebesar $0.3$.
              Berdasarkan nilai parameter terbaik diperoleh hasil $G_{best}$ konvergen pada iterasi ke-$226$ sampai dengan iterasi ke-$900$ dengan nilai fungsi \textit{fitness} konvergen yang dicapai sebesar $0,07442$.

        \item Hasil optimalisasi pemilihan urutan syuting \textit{scene} film "Ketika Adzan Sudah Tidak Lagi Berkumandang" menggunakan algoritma (PSO) memperoleh urutan syuting \textit{scene} yang optimal.
              Urutan syuting \textit{scene} yang diperoleh dimulai dari \textit{scene}
              4-1C-22B-27-10-12B-32A-23B-22C-2A-34B-23A-33B-33A-19-31-1E-7A-32B-24-16B-22A-3-35B-14B-9A-28-12A-29A-21B-35A-37-6A-25-8-\newline
              7B-30A-26-16A-11-2B-19B-1D-13-5-34-6B-20-9B-15-17-16D-18-29B-30B-\newline
              16C-30C-14A-21A-36-4. Total biaya syuting yang diperoleh dengan urutan syuting \textit{scene} tersebut sebesar Rp$7.794.359,48$.

    \end{enumerate}

}
\vspace{-3mm}

% Sub Bab
\section{Saran}
\vspace{-4mm}
{\frenchspacing
    Adapun saran pada penelitian ini sebagai berikut.

    \begin{enumerate}[align = left, left = 0cm, nolistsep]
        \item Penggunaan konsep \textit{traveling salesman problem} (TSP) dapat memberikan hasil lebih akurat jika variabel yang digunakan lebih kompleks atau lebih spesifik, misalnya mempertimbangkan variabel biaya properti tiap \textit{scene}, waktu perjalanan antar lokasi \textit{scene}, durasi syuting, atau batasan lainnya.
        \item Penggunaan algoritma \textit{particle swarm optimization} (PSO) yang telah dimodifikasi dan terbarukan seperti algoritma \textit{Multi-Objective} PSO dapat memberikan solusi optimal dengan model yang lebih kompleks.
    \end{enumerate}
}
\vspace{-3mm}

\pagebreak